\section{从动态规划到 Q 学习}

Bellman 最优性方程已经说明应该求什么：

\[
Q^*(s,a)
=
\E\left[
R_{t+1}
+\gamma\max_{a'}Q^*(S_{t+1},a')
\mid S_t=s,A_t=a
\right].
\]

动态规划的问题是右端需要完整模型：必须知道所有可能的下一状态及其概率。真实机器人通常只得到一次结果——做了动作，传感器返回一个下一状态与一份奖励。能否用这一个随机样本，替代对所有未来的加权求和？

Q 学习的答案是可以，但只能逐步修正。每次交互都提供一份带噪声的 Bellman 目标；学习率把无数次局部修正平均起来，最终逼近同一个不动点。

\subsection{从 Bellman 残差到时序差分误差}

在时刻 \(t\) 观察转移

\[
(S_t,A_t,R_{t+1},S_{t+1}),
\]

用当前估计构造目标

\[
Y_t
=
R_{t+1}
+\gamma\max_{a'}Q_t(S_{t+1},a').
\]

目标与当前预测之差

\[
\delta_t
=Y_t-Q_t(S_t,A_t)
\]

称为时序差分（TD）误差。Q 学习只更新刚访问的状态动作对：

\[
Q_{t+1}(S_t,A_t)
=
Q_t(S_t,A_t)
+\alpha_t\delta_t,
\]

其他表项保持不变。若 \(\delta_t>0\)，这次结果比预期好，就提高估值；若 \(\delta_t<0\)，就下调。更新形式与随机梯度很像，但目标里含当前 \(Q_t\) 自己，称为 bootstrap：用已有估计生成下一次估计。

条件在 \((S_t,A_t)\) 上取期望时，\(Y_t\) 正是 Bellman 最优性算子作用于 \(Q_t\) 的结果。单次样本虽有噪声，却没有系统性偏离正确期望。Q 学习可看成异步、随机的值迭代。

\subsection{为什么叫离策略学习}

生成数据的行为策略记为 \(\mu\)。它可以为了探索而随机选动作；更新目标却总使用

\[
\max_{a'}Q_t(S_{t+1},a'),
\]

即假设下一步会按当前贪心策略行动。收集数据的策略与被评价、被改进的目标策略不同，所以 Q 学习是离策略（off-policy）方法。

与之对照，SARSA 使用实际执行的下一动作 \(A_{t+1}\)：

\[
Q(S_t,A_t)
\leftarrow
Q(S_t,A_t)
+\alpha_t
\left[
R_{t+1}+\gamma Q(S_{t+1},A_{t+1})-Q(S_t,A_t)
\right].
\]

它评价的是包含探索行为在内的当前策略，属于同策略学习。若探索动作有真实风险，两者可能学出不同结果：Q 学习的目标假设未来会贪心，SARSA 会把未来仍可能随机探索造成的风险计入价值。离策略不是天然更好，它只是允许``用一种策略收数据，学习另一种策略''。

\subsection{探索：没试过的动作没有数据}

若从一开始总选当前 \(\argmax_aQ(s,a)\)，初始估计的偶然偏差可能让某些动作永远不被尝试，错误便无法被纠正。最简单的 \(\varepsilon\)-greedy 策略是：以概率 \(1-\varepsilon\) 选当前贪心动作，以概率 \(\varepsilon\) 随机探索。

探索与利用形成真实冲突。探索会牺牲眼前回报来获取信息，利用则把已有知识变成回报。训练模拟器里可以大量探索；医疗、金融和实体机器人里，坏动作的代价不能用``以后会学会''抵消，需要安全约束、离线数据、保守估计或人工规则。MDP 的奖励若没有把安全红线写对，单纯加大碰撞罚分也未必足以阻止罕见灾难。

理论收敛常要求每个状态动作对被访问无穷多次，同时策略最终趋于贪心。这可由逐渐衰减但不太快的 \(\varepsilon_t\) 实现。固定较大的探索率不会收敛到纯贪心行为；衰减太快又可能永远漏掉关键动作。

\subsection{表格情形何时真的收敛}

对有限 MDP，奖励有界、\(0\le\gamma<1\)，若每个状态动作对被访问无穷多次，学习率对每一表项满足 Robbins--Monro 条件

\[
\sum_t\alpha_t=\infty,
\qquad
\sum_t\alpha_t^2<\infty,
\]

则表格 Q 学习在适当条件下几乎必然收敛到 \(Q^*\)。第一条保证新数据永远有累计影响，第二条保证噪声的累计方差有限。典型选择是某表项第 \(n\) 次被访问时取 \(\alpha_n=1/n\)。

这些条件解释了工程现象。固定学习率不会严格收敛到一个常数，而会持续追踪并在其附近波动；在非平稳环境中这反而有用，因为旧数据必须逐渐遗忘。理论里的衰减步长适合固定世界，工程里的固定步长适合移动世界，选择不同是因为问题不同。

收敛定理也只针对表格：每个 \((s,a)\) 有独立可修改的数值。它没有说有限数据时误差多大，没有说探索是否安全，也没有说把表格换成神经网络后仍然成立。

\subsection{一个更新的完整读法}

假设当前估计 \(Q(s,a)=4\)，执行后得到奖励 \(r=1\)，下一状态所有动作中的最大估值为 \(5\)，取 \(\gamma=0.9\)、\(\alpha=0.2\)。目标为

\[
y=1+0.9\times5=5.5,
\]

TD 误差为 \(\delta=5.5-4=1.5\)，更新后

\[
Q_{\mathrm{new}}(s,a)
=4+0.2\times1.5
=4.3.
\]

不要把 \(5.5\) 当成观测到的真实总回报。观测到的只有即时奖励 \(1\) 与下一状态；后面的 \(4.5\) 来自当前模型自己的估计。bootstrap 提高样本效率，也让误差能够自我传播：若下一状态被严重高估，本状态会沿 Bellman 备份继承这份乐观。

最大值还会产生选择偏差。多个动作估计都带噪声时，取最大值更容易挑中正噪声较大的那个，使 Q 值系统性偏高。Double Q-learning 用一套估计选动作、另一套估计评价动作，减弱``同一份噪声既负责选择又负责打分''的偏差。

\subsection{状态太多：从表格到函数逼近}

图像状态不可能逐像素组合建表。于是用参数函数

\[
Q(s,a;\theta)
\]

共享不同状态之间的统计强度，并最小化平方 TD 误差

\[
\ell(\theta)
=
\frac12
\left(
y-Q(s,a;\theta)
\right)^2,
\qquad
y=r+\gamma\max_{a'}Q(s',a';\theta^-).
\]

参数 \(\theta^-\) 暂时视为生成目标的固定参数。梯度下降使一次经验不只修改一个表项，而会改变许多相似状态的估值；这正是泛化能力，也是相互干扰的来源。

深度 Q 网络常用两项稳定化手段：

\begin{itemize}
\tightlist
\item
  经验回放：把转移存入缓冲区并随机抽样，减弱连续样本的强相关，同时重复利用旧经验；
\item
  目标网络：让 \(\theta^-\) 在若干步内固定，再周期性从在线网络复制，避免预测与目标同时快速移动。
\end{itemize}

它们缓解不稳定，却不是普遍收敛证明。回放数据来自旧策略，目标网络仍由自身估计生成；问题的结构没有消失，只是变化速度被压慢。

\subsection{致命三角：三个好主意为何会一起出事}

函数逼近、bootstrap、离策略学习各有明确好处：

\begin{itemize}
\tightlist
\item
  函数逼近让相似状态共享信息；
\item
  bootstrap 不必等整局结束才更新；
\item
  离策略允许复用别的策略、历史系统或人类产生的数据。
\end{itemize}

三者合在一起却可能发散，称为致命三角。直觉是：离策略数据重点覆盖的状态与目标策略不同；函数逼近让一个状态的更新外溢到别处；bootstrap 又把这些被外溢改变的估计当成下一轮目标。误差沿自己生成的数据分布之外循环放大，Bellman 算子在表格 sup 范数下的压缩保证不再自动适用。

线性函数逼近也存在离策略 TD 发散的反例，所以问题不是神经网络太复杂。实际方法会使用目标网络、双重估计、梯度裁剪、保守价值估计或改用策略梯度，但没有一项单独把所有困难清除。

\subsection{离线数据与分布外动作}

若只能使用固定日志数据，学习者不能主动试错。Bellman 目标中的 \(\max_{a'}Q(s',a')\) 可能选中数据里从未出现的动作；函数逼近器对这些动作的高值只是外推，却会被最大值当成事实继续向前传播。这是离线强化学习的核心危险。

行为克隆只模仿数据中的动作，较保守但不会主动超越示范；离线 Q 学习若想改进策略，必须约束新策略别离数据支持太远，或对分布外动作给悲观价值。数据量很大也不保证覆盖了关键反事实：日志只记录当时策略做过什么，没有记录没做的动作会怎样。

\subsection{落点：同一个不动点，不同的信息条件}

把两个单元串起来，结构已经清楚：

\begin{itemize}
\tightlist
\item
  HJB 用已知连续动力学写出无穷小 Bellman 方程；
\item
  离散动态规划用已知转移概率做精确 Bellman 备份；
\item
  Q 学习用样本转移做随机 Bellman 备份；
\item
  函数逼近再把无法存储的值函数压进有限参数。
\end{itemize}

每向现实靠近一步，就少一项信息、多一项近似，也多一种失效方式。强化学习并不保证``只要试得够久就会智能''：状态要近似 Markov，奖励要表达真实目标，探索要覆盖关键动作，数据分布要支撑待评估策略，函数逼近与 bootstrap 还要保持稳定。Bellman 方程给出方向，可信的学习结果则取决于这些前提是否真的成立。
